% !TeX program = pdflatex
\documentclass[11pt,a4paper]{article}

% --- Język i czcionki ---
\usepackage[T1]{fontenc}
\usepackage[utf8]{inputenc}
\usepackage[polish]{babel}
\usepackage{lmodern}
\usepackage[final]{microtype}

% --- Układ strony i nagłówki/stopki ---
\usepackage[a4paper,margin=2.5cm]{geometry}
\usepackage{fancyhdr}
\pagestyle{fancy}
\fancyhf{} % czyść wszystko
\lhead{Etap 2.}
\rhead{\leftmark}
\cfoot{\thepage}

% --- Matematyka ---
\usepackage{amsmath,amssymb,amsthm,mathtools}
\usepackage{siunitx} % jeżeli potrzebne jednostki
\sisetup{locale = PL} % separator 1,23 (polski), można zmienić na locale=PL gdy dostępne
\usepackage{bm}       % pogrubione symbole

% --- Odsyłacze i tytuły ---
\usepackage[hidelinks]{hyperref}
\usepackage[nameinlink,capitalise,noabbrev]{cleveref}
\usepackage{csquotes}
\usepackage{float}

% --- Listy i drobiazgi ---
\usepackage{enumitem}
\setlist{noitemsep,topsep=3pt}
\usepackage{xcolor}

% --- Numeracja i wygląd sekcji ---
\numberwithin{equation}{section}
\setcounter{secnumdepth}{3}
\setcounter{tocdepth}{2}

% --- Meta ---
\title{Sprawozdanie z etapu 3.}
\author{Szymon Hładyszewski, Wiktor Koczkodaj, Jakub Kogut, Wojciech Typer}
\date{\today}

\begin{document}
\maketitle

\section{Raport z testów}
Testy uruchomiono poleceniem \texttt{python manage.py test}. Wykonywane były testy automatyczne aplikacji (testy jednostkowe i integracyjne) zdefiniowane w projekcie. Oczekiwanym wynikiem było pomyślne zakończenie całego zestawu testów bez błędów i niepowodzeń.

\subsection{Zakres testów}
Projekt zawiera kompleksowy zestaw testów dla wszystkich głównych modułów aplikacji:

\begin{itemize}
    \item \textbf{Moduł użytkowników (users):} Testy rejestracji, logowania, zarządzania profilami użytkowników, systemem odznak (badges) oraz awatarów.
    
    \item \textbf{Moduł kursów (courses):} Testy operacji CRUD dla kursów i lekcji, śledzenia postępu użytkowników w lekcjach, uprawnień dostępu do treści.
    
    \item \textbf{Moduł quizów (quizzes):} Testy tworzenia quizów, pytań i odpowiedzi, mechanizmu sprawdzania poprawności odpowiedzi, zapisywania wyników.
    
    \item \textbf{Moduł forum (forum):} Testy tworzenia postów i komentarzy, systemu tagów, uprawnień do edycji i usuwania treści, mechanizmu wyszukiwania.
    
    \item \textbf{Moduł stron (pages):} Testy renderowania szablonów i widoków publicznych.
\end{itemize}

\subsection{Wyniki testów}
Wszystkie testy zakończyły się \textbf{pomyślnie}. System automatycznych testów potwierdził poprawne działanie:
\begin{itemize}
    \item Endpointów API zgodnie ze specyfikacją REST,
    \item Mechanizmów autoryzacji i uwierzytelniania,
    \item Walidacji danych wejściowych,
    \item Relacji między modelami bazy danych,
    \item Logiki biznesowej w poszczególnych modułach,
    \item Migracji bazy danych.
\end{itemize}

Brak błędów w testach potwierdza stabilność systemu oraz gotowość do dalszego rozwoju i wdrożenia.

\end{document}
